% TOP_PROBLEM: {"problemId": "problem.uct-problem-for-nuclear-c-star-algebras", "canonicalTitle": "UCT Problem for Separable Nuclear Complex C*-Algebras", "releaseRank": 203, "primaryDomain": "algebra_representation_category", "primaryDomainLabel": "Algebra, representation & category theory", "displayStatus": "open", "releaseStatus": "open"}
% !TeX program = lualatex
% Definition and sources only; this file is not an LLM solution attempt.
\documentclass[11pt]{article}
\usepackage[margin=25mm]{geometry}
\usepackage{fontspec}
\setmainfont{Latin Modern Roman}
\usepackage{amsmath,amssymb,mathtools}
\usepackage{unicode-math}
\setmathfont{Latin Modern Math}
\usepackage{xurl,hyperref}
\hypersetup{colorlinks=true,urlcolor=blue}
\setcounter{secnumdepth}{0}
\setlength{\parindent}{0pt}
\setlength{\parskip}{5pt}
\setlength{\emergencystretch}{3em}
\begin{document}
\section*{\texorpdfstring{UCT Problem for Separable Nuclear Complex C*-Algebras}{UCT Problem for Separable Nuclear Complex C*-Algebras}}\label{203}
\subsection{Definitions and mathematical statement}
For separable complex $C^*$-algebras $A,B$, let $K_i$ be operator $K$-theory and $KK_j(A,B)$ Kasparov's bivariant groups, graded modulo two. Nuclearity of $A$ means its identity admits point-norm approximations by completely positive contractions factoring through finite-dimensional matrix algebras. Define
\[
 H_j=\bigoplus_{i=0,1}\operatorname{Hom}_{{\mathbb{Z}}}(K_i(A),K_{i+j}(B)),\quad
 E_j=\bigoplus_{i=0,1}\operatorname{Ext}^1_{{\mathbb{Z}}}(K_i(A),K_{i+j+1}(B)).
\]
The problem asks whether every separable nuclear $A$ has the canonical exact sequence
$0\to E_j\to KK_j(A,B)\to H_j\to0$ for every separable $B$ and each $j$. The maps are the Rosenberg--Schochet maps, including the canonical kernel identification; an abstract short exact sequence with arbitrarily chosen maps is insufficient. Neither unitality nor a natural splitting is required.
\par\smallskip\textit{Formulation and concept references:} \href{https://arxiv.org/pdf/1511.02697v3}{[S1]}.

\subsection{Short English statement}
Does operator K-theory, together with its extension term, compute KK-theory through the canonical universal coefficient sequence for every separable nuclear complex C*-algebra?
\subsection{Sources}
\begin{itemize}
\item[S1]Selcuk Barlak and Xin Li, Cartan Subalgebras and the UCT Problem, arXiv1511.02697v3,21June2017,18pages.
\url{https://arxiv.org/pdf/1511.02697v3}.
\textit{Formulation source.}
\end{itemize}
\end{document}
